\chapter{Pengintegralan}\label{ch:b2:integration}

Jilid Tahun ke-1 membangun integral pada sebuah ruas. Bab ini
memperluasnya ke interval sembarang (yakni \omterm{def:b2:integration:improper}{integral tak wajar}, \omterm{def:b2:metric:complete}{lengkap}
dengan kotak perkakas perbandingannya), lalu menelaah integral yang
\emph{bergantung pada sebuah parameter} --- kekontinuan dan penurunan di
bawah tanda integral --- yang digerakkan teorema kekonvergenan
terdominasi, satu-satunya hasil bab ini yang diterima tanpa bukti. Fungsi
$\Gamma$ menjadi contoh yang berjalan terus, sekaligus gerbang menuju
separuh fungsi khusus dalam matematika.

\section{Integral pada interval sembarang}

\begin{definition}\label{def:b2:integration:improper}
Misalkan $f$ \omterm{def:b2:metric:continuity}{kontinu} sepotong-sepotong pada $\intco{a}{b}$ ($b \in \R$
atau $+\infty$). Integralnya disebut
\emph{konvergen}\index{integral tak wajar} apabila $\lim_{x \to b^-}
\int_a^x f$ ada; limitnya lalu ditulis $\int_a^b f$. (Serupa itu pada
$\intoc{a}{b}$, dan pada $\intoo{a}{b}$ dengan memecahnya di sebuah titik
dalam --- pilihan titiknya tak berpengaruh, menurut Chasles.) Integralnya
disebut \emph{konvergen \omterm{def:b2:series:def}{mutlak}} apabila $\int_a^b \abs f$ konvergen;
dan kekonvergenan \omterm{def:b2:series:def}{mutlak} mengakibatkan kekonvergenan, lewat kriteria Cauchy:
\[
\Bigl| \int_x^{y} f \Bigr| \leq \int_x^{y} \abs f
\]
beserta kelengkapan $\R$ (sebab antiturunannya punya sifat Cauchy).
Secara rinci: misalkan $F(x) = \int_a^x f$ dan $G(x) = \int_a^x \abs
f$. Jika $\int^b\abs f$ konvergen, maka $G$ punya limit di $b^-$, jadi
untuk setiap $\varepsilon > 0$ ada $c < b$ dengan $G(y) - G(x)
\leq \varepsilon$ setiap kali $c \leq x \leq y < b$; lalu ungkapan di
atas memindahkan sifat Cauchy itu ke $F$. Untuk sembarang barisan $x_n
\to b^-$, nilai $F(x_n)$ lalu membentuk barisan Cauchy berisi bilangan
real, yang konvergen berkat kelengkapan, dan menyelang-nyelingkan dua
barisan semacam itu menunjukkan limitnya sama untuk semuanya: jadi $F$
punya limit di $b^-$.
\end{definition}

\begin{theorem}[Kotak perkakas perbandingan positif]\label{thm:b2:integration:comparison}
Untuk $f, g \geq 0$ yang \omterm{def:b2:metric:continuity}{kontinu} sepotong-sepotong pada $\intco{a}{b}$:
\begin{enumerate}
  \item $\int_a^b f$ \omterm{def:b2:integration:improper}{konvergen} bila dan hanya bila antiturunannya $x
        \mapsto \int_a^x f$ terbatas;
  \item jika $f \leq g$: kekonvergenan $\int g$ memaksa kekonvergenan
        $\int f$; sedangkan kedivergenannya berpindah ke arah sebaliknya;
  \item jika $f \sim g$ di $b$: kedua integralnya bersifat sama;
  \item skala acuannya: di $+\infty$, $\int^{\infty}
        \frac{\dd t}{t^\alpha}$ \omterm{def:b2:integration:improper}{konvergen} bila dan hanya bila $\alpha >
        1$, dan $\int^\infty \frac{\dd t}{t(\ln t)^\beta}$ bila dan
        hanya bila $\beta > 1$; sedangkan di ujung hingga $b$, integral
        $\int^b \frac{\dd t}{(b - t)^\alpha}$ \omterm{def:b2:integration:improper}{konvergen} bila dan hanya
        bila $\alpha < 1$.
\end{enumerate}
\end{theorem}

\begin{proof}
(1) Antiturunan $F(x) = \int_a^x f$ tidak turun (sebab $f \geq
0$). Jika ia terbatas, maka $\ell = \sup_{x < b}F$ hingga dan
$F(x) \to \ell$: sebab diberikan $\varepsilon > 0$, ada $F(x_0) > \ell -
\varepsilon$, lalu kemonotonannya menjebak $F(x) \in
\intoc{\ell - \varepsilon}{\ell}$ untuk $x_0 \leq x < b$. Jika ia tak
terbatas, maka $F \to +\infty$: jadi divergen.

(2) Dari $f \leq g$: $\int_a^x f \leq \int_a^x g$ untuk setiap $x$;
jika $\int^b g$ \omterm{def:b2:integration:improper}{konvergen}, maka ruas kanannya terbatas, sehingga ruas
kirinya pun terbatas, lalu (1) merampungkannya. Kontraposisinya
memindahkan kedivergenannya ke arah sebaliknya.

(3) Dari $f \sim g$ di $b$ diperoleh $c < b$ dengan
\[
\tfrac12\,g(t) \;\leq\; f(t) \;\leq\; 2\,g(t)
\qquad (c \leq t < b) :
\]
jadi menurut (2) yang diterapkan dua arah pada $\intco{c}{b}$, kedua
integralnya bersifat sama; sedangkan potongan awalnya $\intcc{a}{c}$
merupakan integral wajar dan tidak mengubah apa pun.

(4) Antiturunan gamblangnya: untuk $\alpha \neq 1$ dan $\beta \neq 1$,
\[
\int_c^x \frac{\dd t}{t^\alpha}
= \frac{x^{1-\alpha} - c^{1-\alpha}}{1 - \alpha},
\qquad
\int_c^x \frac{\dd t}{t(\ln t)^\beta}
= \frac{(\ln x)^{1-\beta} - (\ln c)^{1-\beta}}{1 - \beta},
\]
dengan logaritma pada kasus yang dikecualikan: dan itu terbatas ketika
$x \to +\infty$ persis bila $\alpha > 1$, atau $\beta > 1$. Di ujung
yang hingga, penyulihan $u = b - t$ menyusutkannya ke
skala $\int_0 u^{-\alpha}\,\dd u$, yang terbatas bila dan hanya bila
$\alpha < 1$. Terapkan (1) setiap kali.
\end{proof}

\begin{example}[Dua pemanasan, dikerjakan sampai tuntas]\label{ex:b2:integration:warmups}
\emph{(a)} $\displaystyle\int_0^1 \ln t\,\dd t$: integrannya
meledak di $0^+$, tetapi $\abs{\ln t} = o\bigl(t^{-1/2}\bigr)$
di sana (sebab logaritma kalah oleh pangkat), dan $\int_0 t^{-1/2}$
\omterm{def:b2:integration:improper}{konvergen}: jadi \omterm{def:b2:integration:improper}{konvergen} \omterm{def:b2:series:def}{mutlak}. Nilainya, lewat pengintegralan parsial
pada $\intcc{\varepsilon}{1}$:
\[
\int_\varepsilon^1 \ln t\,\dd t
= \bigl[t\ln t - t\bigr]_\varepsilon^1
= -1 - \varepsilon\ln\varepsilon + \varepsilon
\xrightarrow[\varepsilon\to0^+]{} -1 .
\]
\emph{(b)} $\displaystyle\int_0^\infty \frac{\ln t}{1 +
t^2}\,\dd t$: masalahnya ada di kedua ujungnya, jadi pecahlah di $1$.
Di dekat $0$: $\abs{\ln t}$ terintegralkan seperti pada (a); di dekat
$\infty$: $\frac{\ln t}{1+t^2} = o(t^{-3/2})$, jadi \omterm{def:b2:integration:improper}{konvergen} \omterm{def:b2:series:def}{mutlak}.
Penyulihan $t = \frac1u$ memetakan $\intoo{0}{1}$ pada
$\intoo{1}{\infty}$ dan
\[
\int_0^1 \frac{\ln t}{1+t^2}\,\dd t
= \int_1^{\infty} \frac{-\ln u}{1 + u^{-2}}\cdot
\frac{\dd u}{u^2}
= -\int_1^\infty \frac{\ln u}{1+u^2}\,\dd u :
\]
jadi kedua paruhnya saling coret, dan integralnya bernilai $0$.
Pelajaran penutupnya: kesetangkupan terhadap $t \mapsto \frac1t$
berharga satu halaman perhitungan --- dan siasat yang sama sudah
menggerakkan \cref{exo:b2:integration:3}.
\end{example}

\begin{example}[Satu nilai, tiga integral]\label{ex:b2:integration:onecos}
Telaahlah $I = \displaystyle\int_0^{\infty} \frac{1 - \cos
t}{t^2}\,\dd t$. Di $0$: $1 - \cos t \sim \frac{t^2}2$, jadi
integrannya diperluas secara \omterm{def:b2:metric:continuity}{kontinu} oleh nilai $\frac12$ --- sama
sekali tak ada kesingularan. Di $\infty$: $0 \leq \frac{1 - \cos t}{t^2}
\leq \frac{2}{t^2}$, jadi \omterm{def:b2:integration:improper}{konvergen} \omterm{def:b2:series:def}{mutlak}
(\cref{thm:b2:integration:comparison}). Nilainya: integralkan secara
parsial pada $\intcc{\varepsilon}{M}$ dengan $u = 1 - \cos t$ dan $v' =
t^{-2}$:
\[
\int_\varepsilon^M \frac{1 - \cos t}{t^2}\,\dd t
= \Bigl[-\frac{1 - \cos t}{t}\Bigr]_\varepsilon^M
+ \int_\varepsilon^M \frac{\sin t}{t}\,\dd t .
\]
Kurungnya nol di kedua ujungnya (sebab $\frac{1 -
\cos\varepsilon}{\varepsilon} \sim \frac\varepsilon2$, dan pembilangnya
terbatas di $M$), sedangkan integralnya menuju nilai Dirichlet
$\frac\pi2$ (\cref{exo:b2:integration:10}): jadi $I = \frac\pi2$.
Pelajaran penutupnya: dengan $1 - \cos t = 2\sin^2\frac t2$ dan $u =
\frac t2$,
\[
I = \int_0^\infty \frac{2\sin^2 u}{(2u)^2}\;2\,\dd u
= \int_0^\infty \Bigl(\frac{\sin u}{u}\Bigr)^{\!2}\dd u :
\]
ketiga klasiknya, yakni $\int_0^\infty\frac{\sin t}{t}\dd t$,
$\int_0^\infty\bigl(\frac{\sin t}{t}\bigr)^2\dd t$
(\cref{exo:b2:integration:11}) dan $I$, semuanya bernilai
$\frac\pi2$, yang diedarkan lewat pengintegralan parsial dan penyulihan
--- dan hanya yang pertama yang \omterm{def:b2:integration:improper}{semi-konvergen}: sebab pengintegralan
parsialnya menukar kekonvergenan \omterm{def:b2:series:def}{mutlaknya} dengan integran yang lebih sederhana.
\end{example}

\begin{example}[Sebuah integral semi-konvergen]\label{ex:b2:integration:sint}
$\displaystyle\int_1^{\infty} \frac{\sin t}{t}\,\dd t$ \omterm{def:b2:integration:improper}{konvergen}:
integralkan secara parsial,
\[
\int_1^x \frac{\sin t}{t}\dd t
= \Bigl[\frac{-\cos t}{t}\Bigr]_1^x - \int_1^x \frac{\cos
t}{t^2}\dd t ,
\]
dengan kurungnya berlimit dan integral terakhirnya \omterm{def:b2:integration:improper}{konvergen}
\omterm{def:b2:series:def}{mutlak} (sebab $\abs{\cos t}/t^2 \leq t^{-2}$). Namun tidak \omterm{def:b2:series:def}{mutlak}:
sebab dari $\abs{\sin t} \geq \sin^2 t$,
\[
\int_1^x \frac{\abs{\sin t}}{t}\,\dd t
\;\geq\; \int_1^x \frac{\sin^2t}{t}\,\dd t
= \underbrace{\int_1^x \frac{\dd t}{2t}}_{=\ \frac12\ln x
\ \to\ \infty}
\;-\; \underbrace{\int_1^x \frac{\cos 2t}{2t}\,\dd
t}_{\text{konvergen}} ,
\]
dengan integral terakhirnya \omterm{def:b2:integration:improper}{konvergen} lewat pengintegralan parsial yang
\emph{sama} seperti di atas (dengan $\sin 2t$ di kurungnya):
jadi potongan divergen dikurangi potongan \omterm{def:b2:integration:improper}{konvergen} bersifat divergen.
Karenanya $\int_1^\infty\frac{\sin t}{t}\dd t$ \omterm{def:b2:integration:improper}{konvergen} tanpa
\omterm{def:b2:integration:improper}{konvergen} \omterm{def:b2:series:def}{mutlak} --- yakni padanan integral bagi deret
berselang-seling, dengan pengintegralan parsial berperan sebagai
uji berselang-selingnya.
\end{example}

\section{Teorema kekonvergenan}

\begin{theorem}[Kekonvergenan terdominasi]\label{thm:b2:integration:dominated}
Misalkan $(f_n)$ \omterm{def:b2:metric:continuity}{kontinu} sepotong-sepotong pada sebuah interval $I$,
\omterm{def:b2:integration:improper}{konvergen} titik demi titik ke $f$ yang \omterm{def:b2:metric:continuity}{kontinu} sepotong-sepotong, dan
andaikan ada $\varphi \geq 0$ tetap yang terintegralkan ($\int_I \varphi
< \infty$) dengan\index{teorema kekonvergenan terdominasi}
\[
\abs{f_n(t)} \leq \varphi(t) \qquad (t \in I,\ n \in \N).
\]
Maka semua $\int_I f_n$ dan $\int_I f$ \omterm{def:b2:integration:improper}{konvergen} \omterm{def:b2:series:def}{mutlak}, dan
\[
\int_I f_n \xrightarrow[n \to \infty]{} \int_I f .
\]
\end{theorem}
\admitted

\begin{remark}
Bukti jujurnya menjadi milik teori pengintegralan Lebesgue pada Tahun
ke-3; namun pernyataannya dipakai terus-menerus mulai sekarang.
Hipotesis \emph{dominasinya} adalah seluruh intinya: sebab kekonvergenan
titik demi titik saja tidak cukup ($f_n = n\,\mathbf{1}_{\intoo{0}{1/n}}$,
yakni bonggol yang meluncur: $\int f_n = 1 \not\to 0 = \int f$). Teorema
itu juga berlaku untuk parameter yang \omterm{def:b2:metric:continuity}{kontinu} ($f_\lambda$, $\lambda \to
\lambda_0$), lewat pencirian limit secara barisan.
\end{remark}

\begin{example}[Sebuah limit Gauss, lewat dominasi]\label{ex:b2:integration:wallislimit}
Hitunglah $\displaystyle\lim_{n\to\infty} I_n$ dengan $I_n =
\int_0^\infty \Bigl(1 + \frac{t^2}{n}\Bigr)^{\!-n}\dd t$.
Titik demi titik, $(1 + t^2/n)^n \to \eu^{t^2}$ (lewat limit
bunga majemuk), jadi integrannya menuju $\eu^{-t^2}$. Dominasinya:
barisan $n \mapsto (1 + u/n)^n$ tidak turun untuk $u \geq 0$
(sebab AM--GM pada $n + 1$ faktor $1, 1 + \frac un, \dots, 1 +
\frac un$ memberi $(1 + \frac u{n+1})^{n+1} \geq (1 + \frac
un)^n$), jadi untuk $n \geq 2$:
\[
\Bigl(1 + \frac{t^2}{n}\Bigr)^{\!-n}
\leq \Bigl(1 + \frac{t^2}{2}\Bigr)^{\!-2},
\]
yakni pendominasi yang terintegralkan ($\sim 4t^{-4}$ di tak hingga).
Kekonvergenan terdominasi memberi
\[
I_n \xrightarrow[n\to\infty]{} \int_0^\infty \eu^{-t^2}\dd t
= \frac{\sqrt\pi}{2}
\]
(yakni integral Gauss pada \cref{exo:b2:integration:8}). Pemeriksaan
penutupnya: penyulihan $t = \sqrt n\tan\theta$ menghitung $I_n$
secara eksak, $I_n = \sqrt n\int_0^{\pi/2}\cos^{2n-2}\theta\,\dd\theta
= \sqrt n\,W_{2n-2}$, lalu asimtotik Wallis $W_m \sim
\sqrt{\pi/(2m)}$ (\cref{lem:b2:comparison:wallis}) memberi $\sqrt
n\,W_{2n-2} \to \frac{\sqrt\pi}2$ sekali lagi: jadi kedua pilar bab ini
dan bab sebelumnya sepakat.
\end{example}

\begin{example}[Kekonvergenan terdominasi, parameter kontinu]
\label{ex:b2:integration:arctanlimit}
Hitunglah
\[
\lim_{x\to+\infty}\int_0^\infty
\frac{\arctan(xt)}{1+t^2}\,\dd t .
\]
Untuk tiap $t > 0$ berlaku
$\arctan(xt) \to \frac\pi2$ ketika $x \to \infty$; sedangkan
dominasinya,
\[
\Bigl|\frac{\arctan(xt)}{1+t^2}\Bigr|
\leq \frac{\pi/2}{1+t^2},
\qquad\text{terintegralkan, tak bergantung pada } x,
\]
berlaku untuk setiap $x$. Menurut bentuk berparameter \omterm{def:b2:metric:continuity}{kontinu} pada
\cref{thm:b2:integration:dominated} (lewat pencirian secara barisan:
ujilah sepanjang setiap $x_n \to \infty$),
\[
\int_0^\infty\frac{\arctan(xt)}{1+t^2}\,\dd t
\xrightarrow[x\to+\infty]{}
\frac\pi2\int_0^\infty\frac{\dd t}{1+t^2}
= \frac{\pi^2}{4} .
\]
Pelajaran penutupnya: satu titik $t = 0$, tempat limit titik demi
titiknya $0$ dan bukan $\frac\pi2$, tidak mengubah apa pun --- sebab
fungsi limitnya hanya masuk lewat integralnya, salah satu belas kasih
teorema itu yang sunyi.
\end{example}

\section{Integral dengan parameter}

\begin{theorem}[Kekontinuan di bawah tanda integral]\label{thm:b2:integration:continuity}
Misalkan $f \colon A \times I \to \R$ ($A$ sebuah \omterm{def:b2:metric:def}{ruang metrik}, $I$
sebuah interval) dengan: $t \mapsto f(x, t)$ \omterm{def:b2:metric:continuity}{kontinu} sepotong-sepotong
untuk tiap $x$; $x \mapsto f(x, t)$ \omterm{def:b2:metric:continuity}{kontinu} untuk tiap $t$; dan sebuah
\emph{dominasi} $\abs{f(x,t)} \leq \varphi(t)$ (dengan $\varphi$
terintegralkan pada $I$ dan tak bergantung pada $x$). Maka
\[
F(x) = \int_I f(x, t)\,\dd t
\]
terdefinisi dan \omterm{def:b2:metric:continuity}{kontinu} pada $A$.\index{integral berparameter}
\end{theorem}

\begin{proof}
Keterdefinisiannya: dominasinya memberi kekonvergenan \omterm{def:b2:series:def}{mutlak}.
Kekontinuan di $x_0$: untuk sembarang barisan $x_n \to x_0$, fungsi
$g_n(t) = f(x_n, t)$ \omterm{def:b2:integration:improper}{konvergen} titik demi titik ke $f(x_0, t)$ (berkat
kekontinuan pada $x$) di bawah dominasi tetap $\varphi$: jadi
kekonvergenan terdominasi memberi $F(x_n) \to F(x_0)$; lalu rampungkan
lewat pencirian kekontinuan secara barisan (\cref{def:b2:metric:continuity}).
\end{proof}

\begin{theorem}[Penurunan di bawah tanda integral]\label{thm:b2:integration:leibnizrule}
Misalkan $f \colon J \times I \to \R$ ($J$ sebuah interval parameter)
dengan: $t \mapsto f(x,t)$ terintegralkan pada $I$ untuk tiap $x$; $x
\mapsto f(x,t)$ berkelas $C^1$ untuk tiap $t$, dengan turunan parsial
$\frac{\partial f}{\partial x}$ yang \omterm{def:b2:metric:continuity}{kontinu} sepotong-sepotong pada $t$
dan terdominasi: $\bigl|\frac{\partial f}{\partial x}(x,t)\bigr| \leq
\psi(t)$ dengan $\psi$ terintegralkan. Maka $F(x) = \int_I f(x,t)\dd t$
berkelas $C^1$ pada $J$ dan
\[
F'(x) = \int_I \frac{\partial f}{\partial x}(x, t)\,\dd t .
\]
\end{theorem}

\begin{proof}
Tetapkan $x$ dan $h_n \to 0$. Hasil bagi selisihnya
\[
\frac{F(x + h_n) - F(x)}{h_n}
= \int_I \frac{f(x + h_n, t) - f(x, t)}{h_n}\,\dd t
\]
punya integran yang \omterm{def:b2:integration:improper}{konvergen} titik demi titik ke $\frac{\partial
f}{\partial x}(x, t)$, dan terdominasi oleh $\psi(t)$: sebab menurut
ketaksamaan nilai rata-rata yang diterapkan pada $x$ dengan $t$ tetap,
\[
\Bigl|\frac{f(x + h_n, t) - f(x,t)}{h_n}\Bigr|
\leq \sup_{\xi} \Bigl|\frac{\partial f}{\partial x}(\xi, t)\Bigr|
\leq \psi(t) .
\]
Kekonvergenan terdominasi memberi limit $\int_I \frac{\partial
f}{\partial x}(x,t)\dd t$ bagi hasil baginya: jadi $F$ dapat diturunkan
dengan turunan yang diumumkan tadi, yang \omterm{def:b2:metric:continuity}{kontinu} menurut
\cref{thm:b2:integration:continuity} yang diterapkan pada
$\frac{\partial f}{\partial x}$.
\end{proof}

\begin{example}[Sebuah integral berparameter diperiksa terhadap rumusnya]\label{ex:b2:integration:arctanparam}
Misalkan $F(x) = \displaystyle\int_0^\infty \frac{\dd t}{t^2 + x}$
untuk $x > 0$. Pada tiap $\intcc{a}{b} \subset \intoo{0}{\infty}$,
integrannya terdominasi oleh $\frac{1}{t^2 + a}$, yang terintegralkan
dan tak bergantung pada $x$: jadi $F$ \omterm{def:b2:metric:continuity}{kontinu}
(\cref{thm:b2:integration:continuity}). Di sini teoremanya dapat
diperiksa terhadap nilai gamblangnya:
\[
F(x) = \Bigl[\frac{1}{\sqrt x}\arctan\frac{t}{\sqrt
x}\Bigr]_0^\infty = \frac{\pi}{2\sqrt x} ,
\]
yang jelas \omterm{def:b2:metric:continuity}{kontinu}. Kini turunkan di bawah integralnya: turunan
terhadap $x$, yakni $-\frac{1}{(t^2+x)^2}$, terdominasi pada
$\intcc ab$ oleh $\frac{1}{(t^2+a)^2}$ yang terintegralkan, jadi
\cref{thm:b2:integration:leibnizrule} memberi
\[
F'(x) = -\int_0^\infty \frac{\dd t}{(t^2 + x)^2}
\qquad\text{sedangkan}\qquad
F'(x) = -\frac{\pi}{4}\,x^{-3/2} ,
\]
sehingga kita telah \emph{menghitung} sebuah integral baru secara
cuma-cuma: $\int_0^\infty\frac{\dd t}{(t^2+x)^2} = \frac{\pi}{4x^{3/2}}$.
Pelajaran penutupnya: menurunkan sebuah \omterm{thm:b2:integration:continuity}{integral berparameter} yang sudah
dikenal adalah pabrik rumus baru --- dan mengiterasikannya memberi
$\int_0^\infty\frac{\dd t}{(t^2+1)^n}$ untuk setiap $n$, tanpa satu pun
penyulihan trigonometri.
\end{example}

\begin{method}[Menelaah sebuah integral tak wajar]\label{met:b2:integration:study}
Diberikan $\int_a^b f$:
\begin{enumerate}
  \item Temukan masalahnya: daftarlah ujung (atau titik dalam)
        tempat $f$ tak terbatas atau tempat intervalnya
        tak hingga, lalu pecahlah sehingga tiap potongannya punya
        tepat satu ujung yang bermasalah.
  \item Jika $f$ bertanda tetap di dekat ujung itu, carilah
        yang setara lalu bandingkan dengan skala acuan pada
        \cref{thm:b2:integration:comparison}.
  \item Jika $f$ berayun, ujilah $\abs f$ lebih dulu (yakni
        kekonvergenan \omterm{def:b2:series:def}{mutlaknya}). Jika $\int\abs f$ divergen, cobalah
        pengintegralan parsial untuk menukar ayunannya dengan
        peluruhan, seperti pada \cref{ex:b2:integration:sint};
        sedangkan pembatasan bawah seperti $\abs{\sin t} \geq \sin^2t$
        mendeteksi semi-kekonvergenan yang sejati.
  \item Untuk nilainya, bukan sekadar sifatnya: pakailah
        pengintegralan parsial, penyulihan, atau sebuah parameter
        (turunkan integral yang lebih sederhana, seperti pada
        \cref{ex:b2:integration:arctanparam} dan
        \cref{ex:b2:integration:laplace}).
  \item Pemeriksaan kewarasan atas nilai yang terhitung: periksa
        tanda dan ukuran kasarnya terhadap batas yang kasar ($\int_0^\infty
        \eu^{-t^2}\dd t \in \intoo{0}{1 + \int_1^\infty
        \eu^{-t}}$, jadi $\frac{\sqrt\pi}{2} \approx 0.886$ memang
        masuk akal); serta keselarasan dimensinya di bawah penskalaan
        (sebab $t \mapsto \lambda t$ wajib menskalakan kedua ruasnya
        dengan cara yang sama --- yakni pendeteksi tercepat bagi faktor
        yang hilang).
\end{enumerate}
\end{method}

\begin{remark}[Jebakan yang sering muncul]
Ada tiga galat yang berulang. \emph{(i) Pendominasi yang bergantung
pada parameternya:} dominasi $\abs{f(x,t)} \leq \varphi(t)$ wajib
seragam pada $x$ di himpunan yang ditinjau; lazimnya ia berlaku pada
ruas $\intcc ab$ tetapi tidak secara global --- sebab bagi
$\int_0^\infty\eu^{-xt}\dd t$ tak ada pendominasi terintegralkan yang
sahih untuk setiap $x > 0$, namun mendominasinya pada $x \geq a > 0$
sudah cukup untuk bekerja pada seluruh setengah-garis \omterm{def:b2:metric:topology}{terbukanya}, sebab
kekontinuan dan turunan merupakan gagasan yang lokal. \emph{(ii)
Membandingkan integran bertanda:} kotak perkakas perbandingannya untuk
fungsi \emph{tak negatif}; dari $\abs f \leq g$ dengan $\int g$ yang
divergen kita tak dapat menyimpulkan apa pun ---
$\int_1^\infty\frac{\sin t}t\,\dd t$ \omterm{def:b2:integration:improper}{konvergen} walaupun setiap
pembandingan dengan $\frac1t$ gagal. \emph{(iii) Melupakan separuh
masalahnya:} pada $\intoo{0}{\infty}$ selalulah menelaah kedua ujungnya
secara terpisah; sebab $\int_0^\infty\frac{\dd t}{t}$ divergen di
\emph{keduanya}, dan pemecahan yang tampak \omterm{def:b2:integration:improper}{konvergen} dapat diam-diam
mencoret dua ketakhinggaan. Refleks yang aman adalah daftar periksa pada
\cref{met:b2:integration:study}.
\end{remark}

\begin{example}[Sebuah kasus batas Bertrand, sampai ke angkanya]\label{ex:b2:integration:bertrandboundary}
Skala $\int^\infty\frac{\dd t}{t(\ln t)^\beta}$ pada
\cref{thm:b2:integration:comparison} duduk persis di tepi
skala pangkatnya; kasus batasnya pantas mendapat satu
perhitungan penuh. Untuk $\beta = 2$:
\[
\int_\eu^{\infty}\frac{\dd t}{t(\ln t)^2}
= \Bigl[-\frac{1}{\ln t}\Bigr]_\eu^{\infty}
= 0 - (-1) = 1 ,
\]
yakni integral \omterm{def:b2:integration:improper}{konvergen} dengan nilai eksak yang menyenangkan;
sedangkan untuk $\beta = 1$,
\[
\int_\eu^{x}\frac{\dd t}{t\ln t}
= \bigl[\ln\ln t\bigr]_\eu^{x} = \ln\ln x
\longrightarrow \infty ,
\]
jadi divergen --- tetapi begitu lambat sehingga mencapai $\ln\ln x = 10$
menuntut $x = \eu^{\eu^{10}} \approx 10^{9566}$. Pelajaran
penutupnya: di antara ``setiap pangkat $t^{-1-\varepsilon}$
\omterm{def:b2:integration:improper}{konvergen}'' dan ``$t^{-1}$ divergen'' hidup tangga tak hingga
berisi skala logaritmik, yang masing-masing memperhalus yang
sebelumnya; dan penyulihan $u = \ln t$ meruntuhkan tiap anak tangganya
ke anak tangga sebelumnya, sehingga kriteria Bertrand menggemakan
kriteria Riemann satu tingkat di atasnya.
\end{example}

\begin{remark}[Pandangan ke depan di dalam jilid ini]
Perkakas bab ini sebentar lagi ada di mana-mana. Kekonvergenan
terdominasi adalah mesin di balik identitas hampiran pada
bab berikutnya (yakni kernel yang meluncur, baik Bernstein maupun Fejér);
sedangkan kekontinuan dan penurunan di bawah tanda integral
menghasilkan kalkulus koefisien Fourier pada bab Fourier, tempat setiap
$c_n(f)$ merupakan \omterm{thm:b2:integration:continuity}{integral berparameter} yang menyamar. Fungsi $\Gamma$
kembali dua kali: pada bab integral lipat, tempat
sebuah integral ganda akhirnya membuktikan rumus Beta--Gamma milik Euler
selengkapnya, dan pada bab peluang, tempat integral bertipe $\Gamma$
menormalkan kepadatan bakunya lalu menghitung momennya. Adapun
$\int\frac{\sin t}{t}$ yang \omterm{def:b2:integration:improper}{semi-konvergen} muncul kembali sebagai
konstanta Gibbs pada bab Fourier --- integral yang sama, kini menakar
lonjakan jumlah parsialnya di sebuah lompatan.
\end{remark}

\begin{definition}[Fungsi $\Gamma$]\label{def:b2:integration:gamma}
Untuk $x > 0$:\index{fungsi Gamma}
\[
\Gamma(x) = \int_0^{\infty} t^{x-1}\,\eu^{-t}\,\dd t ,
\]
yang \omterm{def:b2:integration:improper}{konvergen} di kedua ujungnya (sebab $t^{x-1}$ terintegralkan di
$0^+$ untuk $x > 0$; dan peluruhan eksponensialnya di $\infty$).
\end{definition}

\begin{theorem}\label{thm:b2:integration:gammaprops}
Fungsi $\Gamma$ \omterm{def:b2:metric:continuity}{kontinu} pada $\intoo{0}{+\infty}$, memenuhi persamaan
fungsional
\[
\Gamma(x + 1) = x\,\Gamma(x),
\qquad \Gamma(1) = 1,
\qquad\text{sehingga}\qquad \Gamma(n + 1) = n! ,
\]
dan berkelas $C^1$ (bahkan $C^\infty$) dengan $\Gamma'(x) =
\int_0^\infty t^{x-1}\eu^{-t}\ln t\,\dd t$.
\end{theorem}

\begin{proof}
Persamaan fungsionalnya: integralkan secara parsial pada
$\intcc{\varepsilon}{M}$ lalu lepaskan ujungnya: $\int t^{x}\eu^{-t} =
[-t^x\eu^{-t}] + x\int t^{x-1}\eu^{-t}$, dengan suku batasnya lenyap ---
memang $\varepsilon^x\eu^{-\varepsilon} \to 0$ ketika $\varepsilon \to
0^+$ karena $x > 0$, dan $M^x\eu^{-M} \to 0$ ketika $M \to \infty$
karena eksponensial mengalahkan setiap pangkat; sedangkan kedua integral
terpotongnya \omterm{def:b2:integration:improper}{konvergen} ke nilai tak wajarnya berkat kekonvergenan yang
sudah ditegakkan pada \cref{def:b2:integration:gamma}. Lalu $\Gamma(1) =
\int \eu^{-t} = 1$; dan induksi memberi faktorialnya.

Kekontinuannya pada $\intcc{a}{b} \subset \intoo{0}{\infty}$: dominasilah
$t^{x-1}\eu^{-t}$ oleh $\varphi(t) = (t^{a-1} + t^{b-1})\eu^{-t}$, yang
terintegralkan dan tak bergantung pada $x \in \intcc{a}{b}$: jadi
\cref{thm:b2:integration:continuity} berlaku pada tiap ruas semacam itu,
sehingga pada seluruh setengah-garisnya. Keterdiferensialannya: turunan
terhadap $x$, yakni $t^{x-1}\eu^{-t}\ln t$, terdominasi pada
$\intcc{a}{b}$ oleh $(t^{a-1} + t^{b-1})\eu^{-t}\,\abs{\ln t}$, yang
tetap terintegralkan: jadi \cref{thm:b2:integration:leibnizrule};
lalu mengiterasikannya memberi semua turunannya (sebab tiap kali
bertambah satu pangkat $\ln t$, yang tak berbahaya).
\end{proof}

\begin{example}[Faktorial setengah bulat]\label{ex:b2:integration:halffactorial}
Persamaan fungsionalnya beserta $\Gamma\bigl(\frac12\bigr) =
\sqrt\pi$ (yang berjarak satu penyulihan dari
\cref{exo:b2:integration:8}: ambil $t = u^2$ pada integral
pendefinisinya) menghasilkan semua nilai setengah bulatnya:
\[
\Gamma\Bigl(\frac32\Bigr) = \frac12\,\Gamma\Bigl(\frac12\Bigr)
= \frac{\sqrt\pi}{2},
\qquad
\Gamma\Bigl(\frac52\Bigr) = \frac32\cdot\frac{\sqrt\pi}{2}
= \frac{3\sqrt\pi}{4},
\qquad
\Gamma\Bigl(\frac72\Bigr) = \frac{15\sqrt\pi}{8} .
\]
Karena $\Gamma(n+1) = n!$, adil kiranya mengatakan ``$\frac12! =
\frac{\sqrt\pi}{2} \approx 0.886$'': jadi faktorialnya telah
diinterpolasi, dan kurva penginterpolasinya menukik di bawah $1$
di antara $0! = 1$ dan $1! = 1$ (dengan minimumnya $\approx 0.8856$ di
$x \approx 1.4616$, yang cocok dengan gambaran kecembungan pada
Bagian I soal akhir pekan). Pelajaran penutupnya: tak ada yang
mengistimewakan bilangan bulat di dalam integral
$\int_0^\infty t^{x-1}\eu^{-t}\dd t$ --- kediskretan faktorialnya
hanyalah kebetulan dari pencacahan, dan $\sqrt\pi$ itulah yang hidup di
antara $1$ dan $1$.
\end{example}

\begin{remark}[Ke mana $\Gamma$ melangkah setelah ini]
Soal akhir pekan bab ini membangun seluruh kalkulus Euler di sekitar
$\Gamma$: yakni \omterm{pb:b2:integration:1}{fungsi Beta}, rekursi pengintegralan parsialnya,
integral Wallis sebagai nilai Beta, dan \omterm{pb:b2:integration:1}{rumus limit Gauss}. Bab
integral lipat membuktikan rumus Beta--Gamma milik Euler untuk semua
argumennya lewat sebuah integral ganda; sedangkan bab peluang menjumpai
$\Gamma$ lagi pada penormalan kepadatan yang paling lazim dan pada momen
waktu tunggu. Jilid Tahun ke-3 membangun ulang $\Gamma$ di atas landasan
Lebesgue, membuktikan teorema ketunggalan Bohr--Mollerup, lalu
memperluas rumus Stirling dari bilangan bulat ke setengah-garis real
lewat kekonvergenan terdominasi.
\end{remark}

\begin{example}[Sebuah perhitungan klasik lewat penurunan]\label{ex:b2:integration:laplace}
Untuk $x \in \R$, misalkan $F(x) = \int_0^{\infty}
\eu^{-t^2}\cos(xt)\,\dd t$ (yang \omterm{def:b2:integration:improper}{konvergen} \omterm{def:b2:series:def}{mutlak}, terdominasi oleh
$\eu^{-t^2}$). Menurut \cref{thm:b2:integration:leibnizrule} (dengan
dominasi turunan terhadap $x$ oleh $t\,\eu^{-t^2}$ yang terintegralkan):
\[
F'(x) = -\int_0^\infty t\,\eu^{-t^2}\sin(xt)\,\dd t
= \Bigl[\frac{\eu^{-t^2}}{2}\sin(xt)\Bigr]_0^\infty
- \frac x2\int_0^\infty \eu^{-t^2}\cos(xt)\,\dd t
= -\frac x2\,F(x),
\]
(lewat pengintegralan parsial dengan $u' = t\eu^{-t^2}$). Persamaan
diferensial $F' = -\frac x2 F$ terintegralkan menjadi $F(x) =
F(0)\,\eu^{-x^2/4}$: jadi integral bertipe Gauss memproduksi dirinya
sendiri. Konstantanya $F(0) = \int_0^\infty \eu^{-t^2}\dd t =
\frac{\sqrt\pi}{2}$ dihitung pada \cref{exo:b2:integration:8} --- dan
sekali lagi, lewat pengintegralan ganda, pada \cref{ch:b2:multint}.
\end{example}

\section{Latihan}

\begin{exercise}[$\star$]\label{exo:b2:integration:1}
Sifat dari:
$\displaystyle\int_0^1 \frac{\dd t}{\sqrt{t(1-t)}}$;
$\;\displaystyle\int_1^\infty \frac{\ln t}{t^2}\dd t$;
$\;\displaystyle\int_0^\infty \frac{\dd t}{1 + t^2\sin^2 t}$
\emph{(bandingkan dengan perilaku bertipe harmonik yang divergen di
dekat $t = n\pi$)}.
\end{exercise}

\begin{exercise}[$\star$]\label{exo:b2:integration:2}
Hitunglah $\displaystyle\int_0^\infty t^n \eu^{-\lambda t}\,\dd t$
($\lambda > 0$) lewat $\Gamma$, dan
$\displaystyle\int_0^1 (\ln t)^n \dd t$ lewat penyulihan $t =
\eu^{-u}$.
\end{exercise}

\begin{exercise}[$\star$]\label{exo:b2:integration:3}
Buktikan bahwa $\displaystyle\int_0^{\infty} \frac{\dd t}{(1 +
t^2)(1 + t^x)}$ terdefinisi dengan baik untuk setiap $x \in \R$ dan
tak bergantung pada $x$. \emph{(Sulihkan $t \mapsto \frac1t$ lalu
rata-ratakan kedua ungkapannya.)} Berapa nilainya?
\end{exercise}

\begin{exercise}[$\star\star$]\label{exo:b2:integration:4}
(Integral Bertrand di ujung yang hingga) Untuk $(\alpha,
\beta)$ yang mana $\displaystyle\int_0^{1/2}
\frac{\dd t}{t^\alpha\,\abs{\ln t}^\beta}$ \omterm{def:b2:integration:improper}{konvergen}?
\end{exercise}

\begin{exercise}[$\star\star$]\label{exo:b2:integration:5}
Misalkan $F(x) = \displaystyle\int_0^{\infty} \frac{\eu^{-xt}}{1 +
t^2}\,\dd t$ untuk $x \geq 0$. Buktikan bahwa $F$ \omterm{def:b2:metric:continuity}{kontinu} pada
$\intco{0}{\infty}$, berkelas $C^2$ pada $\intoo{0}{\infty}$, memenuhi
$F'' + F = \frac1x$ di sana, dan bahwa $F(x) \to 0$ ketika $x \to +\infty$.
\end{exercise}

\begin{exercise}[$\star\star$]\label{exo:b2:integration:6}
(Frullani) Misalkan $f$ \omterm{def:b2:metric:continuity}{kontinu} pada $\intco{0}{+\infty}$ dengan
limit hingga $f(\infty)$ di $+\infty$. Buktikan bahwa untuk $a, b > 0$:
\[
\int_0^{\infty} \frac{f(at) - f(bt)}{t}\,\dd t
= \bigl(f(0) - f(\infty)\bigr)\,\ln\frac ba .
\]
\emph{(Pada $\intcc{\varepsilon}{M}$, sulihkan pada tiap potongannya lalu
kelompokkan ulang menjadi $\int_{a\varepsilon}^{b\varepsilon} -
\int_{aM}^{bM}$ atas $\frac{f(u)}u\,\dd u$; lalu apitlah dengan memakai
kekontinuannya di $0$ dan limitnya di $\infty$.)} Hitunglah
$\int_0^\infty \frac{\eu^{-t} - \eu^{-2t}}{t}\dd t$.
\end{exercise}

\begin{exercise}[$\star\star$]\label{exo:b2:integration:7}
Berilah pembenaran lalu hitunglah $\lim_{n\to\infty}
\displaystyle\int_0^n \Bigl(1 - \frac tn\Bigr)^{\!n} t^{x-1}\,\dd t$
untuk $x > 0$ \emph{(lewat kekonvergenan terdominasi dengan $\varphi(t)
= \eu^{-t}t^{x-1}$, memakai $(1 - t/n)^n \leq \eu^{-t}$; limitnya
$\Gamma(x)$)}.
\end{exercise}

\begin{exercise}[$\star\star\star$]\label{exo:b2:integration:8}
(Integral Gauss lewat siasat parameter) Untuk $x \geq 0$ tetapkan
\[
G(x) = \Bigl(\int_0^x \eu^{-t^2}\dd t\Bigr)^{\!2},
\qquad
H(x) = \int_0^1 \frac{\eu^{-x^2(1+t^2)}}{1 + t^2}\,\dd t .
\]
Buktikan bahwa $G' + H' = 0$ (turunkan $H$ di bawah integralnya lalu
sulihkan $u = xt$ pada integral yang dihasilkannya), lalu turunkan $G(x)
+ H(x) = \frac\pi4$ untuk setiap $x$, dan simpulkan
\[
\int_0^{\infty} \eu^{-t^2}\,\dd t = \frac{\sqrt\pi}{2} .
\]
\end{exercise}

\begin{exercise}[$\star\star\star$]\label{exo:b2:integration:9}
Buktikan bahwa $\Gamma$ bersifat \emph{log-cembung}: yakni $\ln\Gamma$
cembung pada $\intoo{0}{\infty}$. \emph{(Cauchy--Schwarz untuk integral
yang diterapkan pada $t^{(x+y)/2 - 1}\eu^{-t} =
\bigl(t^{x-1}\eu^{-t}\bigr)^{1/2} \bigl(t^{y-1}\eu^{-t}\bigr)^{1/2}$
memberi $\Gamma\bigl(\frac{x+y}{2}\bigr)^2 \leq \Gamma(x)\Gamma(y)$;
lalu padukan dengan kekontinuan dan \cref{exo:b2:realfun:8}.)}
\end{exercise}

\begin{exercise}[$\star\star\star$]\label{exo:b2:integration:10}
(Integral Dirichlet) Tetapkan $F(x) = \displaystyle\int_0^{\infty}
\frac{\sin t}{t}\,\eu^{-xt}\,\dd t$ untuk $x > 0$.
\begin{enumerate}
  \item Berilah pembenaran bagi $F'(x) = -\frac{1}{1 + x^2}$ (turunkan
        di bawah integralnya; lalu hitung $\int_0^\infty
        \eu^{-xt}\sin t\,\dd t$ lewat dua kali pengintegralan parsial).
  \item Buktikan $F(x) \to 0$ ketika $x \to +\infty$ lalu turunkan $F(x)
        = \frac\pi2 - \arctan x$.
  \item Dengan menerima kekontinuan $F$ di $0^+$ (yakni sebuah
        teorema bertipe Abel), simpulkan nilai integral yang
        \omterm{def:b2:integration:improper}{semi-konvergen} itu:
        \[
        \int_0^{\infty} \frac{\sin t}{t}\,\dd t = \frac{\pi}{2}.
        \]
\end{enumerate}
\end{exercise}

\begin{exercise}[$\star\star$]\label{exo:b2:integration:11}
Berilah pembenaran bagi kekonvergenan $\displaystyle\int_0^\infty
\Bigl(\frac{\sin t}{t}\Bigr)^{\!2}\dd t$, lalu hitunglah lewat satu
kali pengintegralan parsial dan \cref{exo:b2:integration:10}:
\[
\int_0^\infty \Bigl(\frac{\sin t}{t}\Bigr)^{\!2}\dd t
= \frac{\pi}{2} .
\]
(Nilainya sama dengan $\int_0^\infty \frac{\sin t}{t}\dd t$ --- tetapi
kali ini kekonvergenannya \omterm{def:b2:series:def}{mutlak}.)
\end{exercise}

\begin{exercise}[$\star\star\star$]\label{exo:b2:integration:12}
(Ekor Gauss) Untuk $x > 0$ tetapkan $T(x) = \displaystyle
\int_x^\infty \eu^{-t^2}\dd t$.
\begin{enumerate}
  \item Dengan menulis $\eu^{-t^2} = \frac{1}{-2t}\cdot(-2t\,\eu^{-t^2})$,
        integralkan secara parsial dua kali untuk memperoleh
        \[
        T(x) = \eu^{-x^2}\Bigl(\frac{1}{2x} -
        \frac{1}{4x^3}\Bigr) +
        \frac34\int_x^\infty \frac{\eu^{-t^2}}{t^4}\,\dd t .
        \]
  \item Batasilah sisanya: $0 \leq \frac34\int_x^\infty
        t^{-4}\eu^{-t^2}\dd t \leq \frac{3}{8x^5}\,\eu^{-x^2}$,
        lalu turunkan pengapitannya
        \[
        \eu^{-x^2}\Bigl(\frac{1}{2x} - \frac{1}{4x^3}\Bigr)
        \leq T(x) \leq \frac{\eu^{-x^2}}{2x},
        \qquad\text{sehingga}\qquad
        T(x) \sim \frac{\eu^{-x^2}}{2x} \quad (x \to +\infty).
        \]
  \item Mengapa deret berselang-seling penuh yang diperoleh dengan
        mengiterasikan pengintegralan parsialnya tak pernah dapat
        \omterm{def:b2:integration:improper}{konvergen} untuk $x$ yang tetap? \emph{(Bandingkan
        pertumbuhan koefisien $1\cdot3\cdots(2k-1)$ dengan
        pangkat $(2x^2)^k$.)}
\end{enumerate}
\end{exercise}

\section{Soal: Integral Euler --- Beta, Gamma, dan rumus limit
Gauss}

\begin{problem}\label{pb:b2:integration:1}
Fungsi $\Gamma$ pada \cref{def:b2:integration:gamma} adalah satu
paruh kalkulus integral milik Euler; paruh satunya adalah
\emph{fungsi Beta}\index{fungsi Beta}
\[
B(x, y) = \int_0^1 t^{x-1}(1 - t)^{y-1}\,\dd t .
\]
Soal ini mengembangkan pasangan $(\Gamma, B)$ hanya dengan perkakas
bab ini --- pengintegralan parsial, penyulihan, dan
kekonvergenan terdominasi --- lalu memuncak pada \emph{rumus
Beta--Gamma milik Euler} $B(x,y) = \frac{\Gamma(x)\Gamma(y)}
{\Gamma(x+y)}$ pada bilangan setengah bulat, dan pada \emph{rumus limit
Gauss}\index{rumus limit Gauss} bagi $\Gamma$. Di sepanjang jalan,
integral Wallis pada \cref{lem:b2:comparison:wallis} muncul kembali
sebagai nilai Beta, dan rumus duplikasi Legendre pun jatuh dengan sendirinya.

\medskip
\noindent\textbf{Bagian I --- Struktur halus $\Gamma$.}
\begin{enumerate}
  \item Ingat kembali mengapa $\Gamma(x) = \int_0^\infty
        t^{x-1}\eu^{-t}\dd t$ \omterm{def:b2:integration:improper}{konvergen} persis untuk $x > 0$, lalu tunjukkan
        \[
        \Gamma(x) \sim \frac1x \qquad (x \to 0^+)
        \]
        \emph{(lewat persamaan fungsionalnya ditambah kekontinuan
        $\Gamma$ di $1$)}.
  \item Buktikan $\Gamma\bigl(\tfrac12\bigr) = \sqrt\pi$
        \emph{(sulihkan $t = u^2$ lalu panggil
        \cref{exo:b2:integration:8})}, lalu turunkan
        $\int_\R \eu^{-u^2/2}\dd u = \sqrt{2\pi}$.
  \item Tunjukkan secara induktif, untuk $n \in \N$:
        \[
        \Gamma\Bigl(n + \frac12\Bigr)
        = \frac{(2n)!}{4^n\,n!}\,\sqrt\pi .
        \]
  \item Berilah pembenaran bagi $\Gamma''(x) = \int_0^\infty
        t^{x-1}\eu^{-t}(\ln t)^2\dd t > 0$, lalu turunkan bahwa
        $\Gamma$ cembung tegas, mencapai minimum tunggal di
        suatu $x_0 \in \intoo{1}{2}$ \emph{(sebab $\Gamma(1) =
        \Gamma(2) = 1$ dan Rolle)}, turun pada $\intoo{0}{x_0}$ lalu
        naik pada $\intoo{x_0}{\infty}$.
  \item Tunjukkan bahwa $\Gamma$ mengalahkan setiap pangkat: untuk tiap
        $k \in \N$ berlaku $x^k = o\bigl(\Gamma(x)\bigr)$ ketika $x \to
        +\infty$ \emph{(apitlah $x$ di antara bilangan bulat lalu pakai
        $\Gamma(n+1) = n!$ beserta kemonotonan pertanyaan 4)}.
\end{enumerate}

\medskip
\noindent\textbf{Bagian II --- \omterm{pb:b2:integration:1}{Fungsi Beta}, lewat pengintegralan parsial.}
\begin{enumerate}[resume]
  \item Tunjukkan bahwa $B(x,y)$ \omterm{def:b2:integration:improper}{konvergen} persis untuk $x > 0$ dan $y
        > 0$, dan bahwa $B(x,y) = B(y,x)$.
  \item Hitunglah $B(x, 1) = \frac1x$, lalu buktikan lewat
        pengintegralan parsial, untuk $x, y > 0$:
        \[
        B(x, y+1) = \frac{y}{x}\,B(x+1, y) .
        \]
  \item Dari pemecahan $t^{x-1}(1-t)^{y-1} =
        t^{x}(1-t)^{y-1} + t^{x-1}(1-t)^{y}$ turunkan $B(x,y) =
        B(x+1,y) + B(x,y+1)$, lalu padukan dengan pertanyaan 7 menjadi
        \emph{relasi penurunan}
        \[
        B(x, y+1) = \frac{y}{x+y}\,B(x,y),
        \qquad
        B(x+1, y) = \frac{x}{x+y}\,B(x,y) .
        \]
  \item Turunkan, untuk bilangan bulat $m, n \geq 1$:
        \[
        B(m, n) = \frac{(m-1)!\,(n-1)!}{(m+n-1)!}
        = \frac{1}{(m+n-1)\binom{m+n-2}{m-1}} .
        \]
  \item Buktikan \emph{rumus Euler untuk satu argumen bulat}:
        untuk setiap $x > 0$ dan $n \in \N^*$,
        \[
        B(x, n) = \frac{\Gamma(x)\,\Gamma(n)}{\Gamma(x + n)}
        \]
        \emph{(secara induktif pada $n$: kedua ruasnya sama dengan
        $\frac1x$ di $n = 1$ dan menuruti relasi penurunan yang sama)}.
\end{enumerate}

\medskip
\noindent\textbf{Bagian III --- Integral Wallis sebagai nilai Beta.}
\begin{enumerate}[resume]
  \item Sulihkan $t = \sin^2\theta$ untuk memperoleh bentuk
        trigonometrinya
        \[
        B(x, y) = 2\int_0^{\pi/2}
        \sin^{2x-1}\theta\,\cos^{2y-1}\theta\,\dd\theta .
        \]
  \item Turunkan $W_n = \frac12\,B\bigl(\frac{n+1}2, \frac12\bigr)$
        bagi integral Wallis $W_n = \int_0^{\pi/2}\sin^n
        \theta\,\dd\theta$, lalu peroleh kembali rekurensi $W_n =
        \frac{n-1}{n}W_{n-2}$ pada \cref{lem:b2:comparison:wallis}
        hanya dari relasi penurunan pertanyaan 8.
  \item Hitunglah $B\bigl(\frac12, \frac12\bigr) = 2W_0 = \pi$ lalu
        periksalah terhadap $\Gamma\bigl(\frac12\bigr)^2/\Gamma(1)$:
        jadi rumus Euler berlaku di $\bigl(\frac12, \frac12\bigr)$.
  \item Turunkan bentuk tertutup $W_{2n} = \frac\pi2\,
        \frac{(2n)!}{4^n(n!)^2}$ dari rekurensinya, lalu periksalah
        \[
        B\Bigl(n + \frac12, \frac12\Bigr)
        = \frac{\Gamma\bigl(n + \frac12\bigr)\Gamma\bigl(
        \frac12\bigr)}{\Gamma(n+1)} .
        \]
        Simpulkan, secara induktif dengan relasi penurunannya, bahwa
        rumus Euler $B(x,y) =
        \frac{\Gamma(x)\Gamma(y)}{\Gamma(x+y)}$ berlaku setiap kali
        $2x$ dan $2y$ merupakan bilangan bulat positif.
  \item Sulihkan $u = \frac{t}{1-t}$ untuk memperoleh bentuk klasik
        ketiganya
        \[
        B(x,y) = \int_0^\infty
        \frac{u^{x-1}}{(1+u)^{x+y}}\,\dd u ,
        \]
        lalu periksalah kasus $x = y = \frac12$ secara langsung (sebab
        $u = v^2$ menyusutkannya menjadi $\int_0^\infty\frac{2\,\dd v}{1+v^2}$).
\end{enumerate}

\medskip
\noindent\textbf{Bagian IV --- \omterm{pb:b2:integration:1}{Rumus limit Gauss}.}
\begin{enumerate}[resume]
  \item Untuk $x > 0$ dan $n \in \N^*$, buktikan lewat $n$ kali
        pengintegralan parsial berturut-turut:
        \[
        \int_0^n \Bigl(1 - \frac tn\Bigr)^{\!n} t^{x-1}\,\dd t
        = \frac{n!\;n^x}{x(x+1)\cdots(x+n)} .
        \]
  \item Rampungkan dengan \cref{exo:b2:integration:7} (yakni
        kekonvergenan terdominasi) sehingga diperoleh \emph{\omterm{pb:b2:integration:1}{rumus limit Gauss}}:
        \[
        \Gamma(x) = \lim_{n\to\infty}
        \frac{n!\;n^x}{x(x+1)\cdots(x+n)}
        \qquad (x > 0).
        \]
  \item Dengan mengambil logaritmanya, tunjukkan bahwa untuk $x > 0$:
        \[
        \ln\Gamma(x) = -\ln x - \gamma x +
        \sum_{k=1}^{\infty}\Bigl(\frac xk -
        \ln\Bigl(1 + \frac xk\Bigr)\Bigr),
        \]
        dengan $\gamma$ menyatakan konstanta Euler
        (\cref{ex:b2:comparison:harmonic}); lalu berilah pembenaran bagi
        kekonvergenan deretnya \emph{(sebab suku umumnya
        $\sim \frac{x^2}{2k^2}$)}.
  \item Pakailah rumus Gauss di $x = \frac12$ beserta asimtotik
        koefisien binomial pusat $\binom{2n}{n} \sim
        \frac{4^n}{\sqrt{\pi n}}$
        (\cref{ex:b2:comparison:centralbinomial}) untuk menghitung ulang
        $\Gamma\bigl(\frac12\bigr) = \sqrt\pi$: jadi konstanta
        Stirling dan integral Gauss adalah bilangan yang sama
        dalam dua samaran.
  \item Periksalah bahwa rumus Gauss membuktikan ulang persamaan
        fungsionalnya: dari kesamaan eksak
        \[
        \frac{n!\,n^{x+1}}{(x+1)\cdots(x+n+1)}
        = \frac{n!\,n^{x}}{x(x+1)\cdots(x+n)}\cdot
        \frac{n\,x}{x+n+1},
        \]
        simpulkan $\Gamma(x+1) = x\,\Gamma(x)$ sekali lagi. (Rumus
        Gauss menentukan $\Gamma$ secara tuntas; sedangkan jilid Tahun
        ke-3 membuktikan teorema Bohr--Mollerup yang lebih tajam:
        yakni persamaan fungsionalnya ditambah kelog-cembungannya sudah
        mengunci $\Gamma$.)
\end{enumerate}

\medskip
\noindent\textbf{Bagian V --- Panennya.}
\begin{enumerate}[resume]
  \item Untuk $a > 0$ tunjukkan $\int_0^\infty \eu^{-t^a}\dd t =
        \Gamma\bigl(1 + \frac1a\bigr)$, lalu hitunglah limitnya ketika
        $a \to +\infty$ lewat kekonvergenan terdominasi \emph{(dengan
        limit titik demi titik $\mathbf 1_{t < 1}$; dan dominasilah oleh
        $1$ pada $\intoc{0}{1}$ serta oleh $\eu^{-t^2}$ di luarnya,
        untuk $a \geq 2$)}. Periksalah jawabannya terhadap kekontinuan
        $\Gamma$.
  \item Untuk $n \geq 1$ tunjukkan
        \[
        \int_0^1 \frac{\dd t}{\sqrt{1 - t^n}}
        = \frac1n\,B\Bigl(\frac1n, \frac12\Bigr),
        \]
        lalu peroleh kembali nilai $2$ (untuk $n = 1$) dan $\frac\pi2$
        (untuk $n = 2$). (Untuk $n = 4$ inilah \emph{konstanta
        lemniskat}, yang tak punya bentuk tertutup yang elementer;
        kisahnya menjadi bagian teori integral eliptik.)
  \item (Momen) Untuk $x > 0$ dan $k \in \N$, tunjukkan
        \[
        \frac{1}{\Gamma(x)}\int_0^\infty
        t^{k}\,t^{x-1}\eu^{-t}\,\dd t
        = \frac{\Gamma(x+k)}{\Gamma(x)}
        = x(x+1)\cdots(x+k-1),
        \]
        yakni faktorial naiknya; lalu periksalah bahwa $x = 1$ memberi
        $k!$. (Pada bab peluang, inilah momen ke-$k$ sebuah kepadatan
        waktu tunggu yang baku.)
  \item Buktikan kesamaan Beta berikut, yang sahih untuk \emph{setiap}
        $x > 0$:
        \[
        B(x, x) = 2^{1-2x}\,B\Bigl(x, \frac12\Bigr)
        \]
        \emph{(sulihkan $t = \frac{1+s}2$, manfaatkan
        kesetangkupannya pada $s$, lalu ambil $s = \sqrt v$)}. Lalu
        turunkan, untuk $2x \in \N^*$, \emph{rumus duplikasi Legendre}
        \[
        \Gamma(x)\,\Gamma\Bigl(x + \frac12\Bigr)
        = 2^{1-2x}\,\sqrt\pi\;\Gamma(2x),
        \]
        dan periksalah secara langsung di $x = n$ lewat pertanyaan 3.
        (Untuk $x$ yang umum ia menyusul dari kesamaan yang sama begitu
        rumus Euler diketahui untuk semua argumennya --- lewat bukti
        integral ganda pada bab integral lipat.)
  \item Rangkuman. Satu kalimat untuk masing-masing: (i) di mana
        pengintegralan parsial mengusung seluruh Bagian II; (ii) di mana
        kekonvergenan terdominasi masuk ke Bagian IV dan V; (iii)
        masukan asimtotik mana yang diimpor dari bab
        perbandingan; (iv) apa yang kini terbukti dari rumus Euler
        $B(x,y) = \Gamma(x)\Gamma(y)/\Gamma(x+y)$, dan apa yang
        tersisa untuk diselesaikan integral gandanya.

\end{enumerate}
\end{problem}
